\bookchapter{Глава 24: Root-права и BL-блокировка: глубокий разбор}

\chapterintro{Эта глава — углублённое расширение главы 20, предназначенное для тех, кто действительно хочет понять низкоуровневые механизмы телефона или даже попробовать прошивку. Обратите внимание: эти операции сопряжены с риском, действуйте осторожно.}

\booksection{Глубокое понимание Bootloader}

Bootloader — первая программа, запускающаяся при включении телефона; её задача — инициализировать оборудование и загрузить ядро операционной системы. Можно считать Bootloader «стражем ворот» телефона.

\booksubsection{Принцип работы BL-блокировки:}

\begin{itemize}
  \item \textbf{Проверка подписи}: BL-блокировка проверяет цифровую подпись загрузочного образа; только система, подписанная производителем, может загрузиться.
  \item \textbf{Защита от понижения версии}: не позволяет прошить старую версию системы, чтобы использовать известные уязвимости.
  \item \textbf{Защита данных}: даже если прошита новая система, без оригинального аккаунта и пароля пользовательские данные недоступны.
\end{itemize}

\booksubsection{Аппаратное хранение и верификация}

Состояние BL-блокировки (Locked/Unlocked) и публичный ключ производителя обычно хранятся в защищённой аппаратной области (например, раздел devinfo у Qualcomm, eFuse или раздел RPMB), что предотвращает программное вмешательство. При каждой загрузке Bootloader считывает состояние из этой безопасной зоны и решает, выполнять ли строгую проверку подписи.

\booksubsection{Проверка цифровой подписи}

При выпуске обновления или образа системы производитель подписывает его своим закрытым ключом. При загрузке устройства Bootloader проверяет системный образ (boot.img, recovery.img) встроенным открытым ключом. Алгоритм:

\begin{enumerate}
  \item Заново вычислить хеш системного образа.
  \item Расшифровать прилагаемую к образу подпись открытым ключом, получив хеш.
  \item Сравнить два хеша; при совпадении проверка пройдена. Это гарантирует, что образ не был подменён третьей стороной.
\end{enumerate}

\booksubsection{Цепочка безопасной загрузки}

BL-блокировка — первое звено цепочки безопасной загрузки. Пройдя проверку, она загружает и запускает следующий этап загрузчика (например, ядро), послойно гарантируя доверенность среды загрузки.

\booksubsection{Блокировка vs разблокировка: двустороннее влияние BL}

Для наглядности покажем различия по 5 ключевым измерениям:

\begin{center}
\begin{tabularx}{\linewidth}{|X|X|X|}
\hline
\textbf{Ключевое измерение} & \textbf{Заблокировано (по умолчанию)} & \textbf{Разблокировано} \tabularnewline
\hline
Безопасность системы & Запрещает прошивку неподписанных неофициальных систем, защищает целостность и низкоуровневую безопасность\newline Высокая безопасность: загружаются только образы с официальной подписью, эффективно против зловредной модификации прошивки & Безопасность снижена: разрешена прошивка сторонних ROM, проверка целостности отключена, зловредное ПО с Root-правами может глубоко модифицировать систему \tabularnewline
Стабильность системы & Гарантирует работу официальной системы, протестированной и оптимизированной производителем\newline Стабильно: официальная прошивка, полный функционал, надёжные OTA-обновления & Риск выше: сторонний ROM может иметь проблемы совместимости — система может крашиться, окирпичиться, аппаратные функции (камера, быстрая зарядка) могут работать некорректно, OTA-обновления могут не работать \tabularnewline
Гарантия & Производитель через заблокированный BL обеспечивает стабильную официальную систему и гарантийное обслуживание\newline Гарантия действительна: бесплатный официальный ремонт (при неумышленных поломках) & Гарантия аннулирована: большинство производителей считают разблокировку BL добровольным отказом от гарантии; даже повторная блокировка может быть обнаружена по истории, и в ремонте откажут \tabularnewline
Безопасность данных & Разблокировка обычно принудительно очищает данные, чтобы предотвратить зловредное использование старых данных в новой системе\newline Данные в безопасности: при повседневном использовании BL ограничивает возможность зловредных программ модифицировать системные разделы & Риск для данных: при разблокировке все пользовательские данные принудительно удаляются; если после разблокировки шифрование настроено неверно, хранилище может быть небезопасно; при потере устройства данные легче извлечь через прошивку \tabularnewline
Функциональные ограничения & Гарантирует нормальную работу функций, зависящих от безопасной среды (платёж, банковские приложения)\newline Полный функционал: все официальные функции работают & Ограниченный функционал: некоторые приложения, чувствительные к безопасности (Google Pay, банковские приложения) могут отказать в работе при обнаружении разблокированного BL; DRM-сервисы могут быть понижены, что влияет на воспроизведение HD-видео \tabularnewline
\hline
\end{tabularx}
\end{center}

\booksection{Подробная процедура разблокировки BL (на примере Xiaomi)}

\begin{enumerate}
  \item \textbf{Подать заявку на разблокировку}: на сайте инструмента разблокировки Xiaomi.
  \item \textbf{Привязать аккаунт}: в настройках телефона привяжите аккаунт Xiaomi.
  \item \textbf{Период ожидания}: в зависимости от уровня аккаунта — от 7 до 30 дней.
  \item \textbf{Разблокировка}: используйте инструмент разблокировки Xiaomi, подключив телефон к компьютеру.
\end{enumerate}

\booksection{Различные решения для Root}

\begin{itemize}
  \item \textbf{Magisk}: современное решение Root с архитектурой Systemless, позволяет скрыть статус Root.
  \item \textbf{Kernelsu}: решение Root на уровне ядра, более низкоуровневое, лучше скрывает наличие Root.
  \item \textbf{SuperSU}: старое решение Root, обновление сейчас прекращено.
\end{itemize}

\booksection{Введение в прошивку}

Прошивка — установка на телефон новой операционной системы. Типы:

\begin{itemize}
  \item \textbf{Прошивка официального ROM}: установка официальной системы, для восстановления или понижения версии.
  \item \textbf{Прошивка стороннего Recovery}: например TWRP, для бэкапа/восстановления и прошивки сторонних ROM.
  \item \textbf{Прошивка стороннего ROM}: например LineageOS, Pixel Experience и др., для другого пользовательского опыта.
\end{itemize}

\booksection{Samsung Knox: необратимая цена разблокировки BL}

В отличие от других Android-производителей, Samsung имеет уникальный дизайн «пережигаемого предохранителя» при разблокировке BL:\textbf{для Samsung разблокировка вызывает пережигание Knox с навсегда утратой некоторых функций безопасности (таких как Samsung Pay); все защищённые данные удаляются; гарантия затрагивается}. Подробнее:

\begin{itemize}
  \item \textbf{Предохранитель одноразовый}: внутри телефона есть eFuse-предохранитель; при обнаружении разблокировки BL он навсегда пережигается — восстановить невозможно.
  \item \textbf{Затронутые функции}: Samsung Pay (транспортные/банковские карты), безопасная папка Knox, корпоративное MDM-управление, некоторые приложения медицинского страхования перестают работать.
  \item \textbf{Принудительная очистка данных}: при срабатывании предохранителя телефон принудительно сбрасывается до заводских настроек; все локальные данные стираются (облачные не затрагиваются).
  \item \textbf{Влияние на гарантию}: производитель прямо заявляет, что после пережигания Knox гарантийное обслуживание не предоставляется.
\end{itemize}

\begin{tcbwarn}{Трижды подумайте}
Пережигание Knox у Samsung\textbf{необратимо}— однажды сработав, вернуть заводское состояние невозможно. Если вам нужны функции, требующие Root, сначала проверьте, есть ли альтернативы без Root.

\end{tcbwarn}

\textbf{Важное предупреждение}\par

Прошивка сопряжена с риском. Возможные последствия:

\begin{itemize}
  \item Окирпичивание телефона (не включается)
  \item Потеря гарантии
  \item Потеря данных
  \item Нарушение работы платёжных функций
\end{itemize}

Перед прошивкой обязательно сделайте полный бэкап, внимательно прочитайте инструкцию и убедитесь в соответствии модели и версии.

\begin{tcbinfo}{Рекомендуемые ресурсы}
\begin{itemize}
  \item \textbf{XDA Developers}: крупнейшее в мире сообщество Android-разработчиков, богатые инструкции по прошивке и ROM-ресурсы.
  \item \textbf{Coolapk}: известное китайское цифровое сообщество, много китайскоязычных инструкций и обсуждений.
  \item \textbf{Форумы конкретных моделей}: сообщество Xiaomi, сообщество OnePlus и др., с официальными и народными инструкциями.
\end{itemize}

\end{tcbinfo}

\booksection{Часто задаваемые вопросы (Q\& A)}

\textbf{Q1: Можно ли восстановить BL-блокировку после разблокировки?}\par

\textbf{A:}Да, можно заблокировать обратно. Обычно нужно прошить официальный полный образ через Fastboot и использовать специальный скрипт (например,\texttt{flash\_all\_lock.\allowbreak bat}у Xiaomi). Однако обратите внимание: повторная блокировка может не обратить некоторые пережигания безопасности (например, Knox у Samsung).

\textbf{Q2: Как получить Root после разблокировки?}\par

\textbf{A:}Разблокировка BL — лишь первый шаг. Root обычно требует:

\begin{enumerate}
  \item Прошивки пользовательского Recovery (например TWRP).
  \item Через Recovery — прошивки Magisk или другого Root-менеджера для патча boot-образа и получения Root-прав.
\end{enumerate}
